\section{La sécurité des technologies de l'information}
\label{la-securite-des-technologies-de-l-information}

\begin{enumerate}
 \item
  Tous les fonctionnaires de la FCEG doivent respecter les politiques suivantes lorsqu\textquotesingle ils utilisent des services électroniques fournis par la FCEG ou des services électroniques utilisés à des fins officielles par la FCEG.
 \item
  Les courriels officiels de la FCEG ou une alternative appropriée déterminée par le/la commissaire à l\textquotesingle informatique doivent être utilisés pour toutes les communications.
 \item
  Tout nom de domaine affilié à la FCEG ou à l\textquotesingle un de ses groupes doit être dûment enregistré au nom de la FCEG.
 \item
  Lorsqu\textquotesingle un officier quitte son poste à la FCEG, les mots de passe de tous les comptes doivent être réinitialisés dans les 48 heures.
 \item
  Le commissaire à l'informatique doit être en mesure de réinitialiser tout mot de passe associé à un compte utilisé à des fins officielles.
 \item
  Toute donnée créée, stockée, transférée, diffusée ou hébergée sur un serveur ou un ordinateur appartenant à la FCEG ou relevant de la compétence du/de la commissaire à l\textquotesingle informatique doit être conforme aux politiques spécifiées dans la politique de traitement et de sécurité de l\textquotesingle information en ce qui concerne le transfert et le stockage de documents classifiés, le cas échéant.
 \item
  Des mots de passe forts doivent être utilisés pour tous les comptes.

  \begin{enumerate}
   \item
    Un mot de passe fort est un mot qui comporte au moins 8 caractères et au moins un de chacun des éléments suivants: lettres majuscules, lettres minuscules, chiffres, caractères spéciaux.
   \item
    Des exceptions peuvent être approuvées par le/la commissaire à l'informatique Ces demandes ne doivent jamais divulguer le mot de passe exact, mais simplement décrire le format du mot de passe et la raison pour laquelle il ne peut pas répondre aux exigences spécifiées.
  \end{enumerate}
 \item
  Les documents classifiés ou autrement privés sont stockés soit dans un système de stockage en nuage approuvé par le commissaire aux technologies de l\textquotesingle information, soit sur un support de stockage crypté.

  \begin{enumerate}
   \item
    Ces documents ne peuvent être consultés que par les fonctionnaires qui ont besoin des informations qu\textquotesingle ils contiennent.
   \item
    Ces documents ne doivent pas être stockés sur des supports personnels, y compris, mais sans s\textquotesingle y limiter, les disques durs des ordinateurs portables, les systèmes de stockage en nuage non approuvés ou les supports de stockage amovibles personnels non cryptés.
  \end{enumerate}
 \item
  Le/la commissaire à l\textquotesingle informatique effectue des tests de pénétration sur l\textquotesingle infrastructure électronique de la FCEG au moins une fois par trimestre.
\end{enumerate}

